\section{Experience Format Specification}

\label{sec:format}

\subsection{Semantic Advantage Definition}

\begin{definition}[Semantic Advantage]
A semantic advantage is a natural language statement $e \in \mathcal{E}$ that encodes a generalizable reasoning pattern extracted from comparative analysis of model trajectories. Formally, $e$ consists of:

\begin{itemize}
    \item \textbf{Text} $e.\text{text} \in \Sigma^*$: The insight itself (max 32 words)
    \item \textbf{Domain} $e.\text{domain} \in \mathcal{D}$: Problem category (e.g., "math.geometry")
    \item \textbf{Confidence} $e.\text{conf} \in [0,1]$: Quality score derived from advantage magnitude
    \item \textbf{Embedding} $e.\text{emb} \in \mathbb{R}^{7680}$: Semantic vector for retrieval
    \item \textbf{Metadata} $e.\text{meta}$: Creation timestamp, author, usage statistics
\end{itemize}
\end{definition}

\subsection{JSON Schema}

The canonical representation uses JSON with content-addressable identifiers:

\begin{lstlisting}[language=json,caption=Experience JSON Schema,label=lst:schema]
{
  "id": "exp_sha256_abc123...",
  "version": "1.0",
  "domain": "math.geometry",
  "text": "When solving geometry problems with 
           intersections, validate solutions lie 
           within bounded regions, not extensions.",
  "confidence": 0.87,
  "examples": [
    {
      "input": "Find intersection of line AB...",
      "output": "Point P lies outside segment AB...",
      "correct": false
    },
    {
      "input": "Find intersection within triangle...",
      "output": "Point Q is at (3,4) inside ABC.",
      "correct": true
    }
  ],
  "metadata": {
    "created_at": "2025-09-15T14:30:00Z",
    "created_by": "0x742d35Cc6634C0532925a...",
    "votes": {"upvotes": 24, "downvotes": 2},
    "usage_count": 156,
    "domains_applied": ["geometry", "optimization"]
  },
  "embedding": [0.123, -0.456, ...],  // 7680-dim
  "merkle_proof": "0xdef456..."
}
\end{lstlisting}

\subsection{Constraints and Validation}

Valid experiences must satisfy:

\begin{enumerate}
    \item \textbf{Length}: $|e.\text{text}| \leq 32$ words (enforced via tokenization)
    \item \textbf{Structure}: "When [context], [action]" or "[precondition] implies [strategy]"
    \item \textbf{Generality}: No problem-specific constants (e.g., "solve $x^2 = 4$" is invalid)
    \item \textbf{Actionability}: Must guide reasoning, not merely state facts
    \item \textbf{Non-Redundancy}: $\forall e' \in \mathcal{E}, \cos(e.\text{emb}, e'.\text{emb}) < 0.9$
\end{enumerate}

\subsection{Example Experience Library}

\begin{example}[Mathematical Reasoning Experiences]
\begin{enumerate}
    \item \texttt{[G0]} "When solving geometry problems with intersections, validate solutions lie within bounded regions or segments, not on extensions, to avoid extraneous answers."
    
    \item \texttt{[G1]} "For expected extreme statistics in combinatorial problems, use direct enumeration for small sizes."
    
    \item \texttt{[G10]} "When using mathematical invariants to prove impossibility, always validate them against known achievable states or small cases."
    
    \item \texttt{[G21]} "For complex polynomials with real parameters, separate real and imaginary parts to find when real roots exist."
\end{enumerate}
\end{example}

These experiences exhibit the desired properties: concise, strategic, domain-agnostic within their category, and immediately applicable.

